%%%%%%%%%%%%%%%%%%%%%%%%%%%%%%%%%%%%%%%%%%%%%%%%%%%%%%%%%%%%
% 12.5 COMPUTABILITY THEORY
% The Composition of Advanced Mathematics
%%%%%%%%%%%%%%%%%%%%%%%%%%%%%%%%%%%%%%%%%%%%%%%%%%%%%%%%%%%%

\section{Computability Theory}
\label{sec:computability-theory}

\prerequisite{
Formal logic, proof theory, functions, sets, relations, countability,
mathematical induction, algorithms, discrete mathematics, and elementary
number theory.
}

\learningobjectives{
By the end of this section, the reader should be able to:
\begin{itemize}
    \item explain the purpose of computability theory;
    \item distinguish informal algorithms from formal models of computation;
    \item understand effective procedures conceptually;
    \item recognize Turing machines;
    \item describe tape, states, symbols, transitions, and configurations;
    \item understand deterministic Turing-machine computation;
    \item recognize halting and nonhalting behavior;
    \item distinguish recognition from decision;
    \item define computable functions;
    \item recognize partial computable functions;
    \item understand semidecidable or recognizable sets;
    \item distinguish decidable and recognizable sets;
    \item understand enumerability;
    \item recognize recursively enumerable sets;
    \item understand enumerators;
    \item recognize equivalent models of computation;
    \item understand the Church--Turing thesis;
    \item recognize register machines and lambda calculus conceptually;
    \item understand universal computation;
    \item recognize universal Turing machines;
    \item understand encoding of machines and inputs;
    \item recognize self-reference in computation;
    \item understand the Halting Problem;
    \item prove the Halting Problem undecidable;
    \item recognize diagonal arguments in computability;
    \item understand mapping reductions;
    \item use reductions to prove undecidability;
    \item recognize many-one reducibility;
    \item understand closure properties of decidable languages;
    \item understand closure properties of recognizable languages;
    \item distinguish a language from its complement;
    \item understand the relationship between recognizability and
          co-recognizability;
    \item recognize Rice's theorem;
    \item understand semantic properties of programs;
    \item recognize undecidability of equivalence problems;
    \item understand Post's correspondence problem conceptually;
    \item recognize undecidable problems in formal systems;
    \item understand computation over strings and natural numbers;
    \item understand pairing and coding functions;
    \item recognize universal computable functions;
    \item understand bounded simulation;
    \item recognize dovetailing;
    \item understand effective enumeration;
    \item distinguish computability from computational complexity;
    \item recognize decidable but computationally difficult problems;
    \item understand oracle computation conceptually;
    \item recognize relative computability;
    \item understand Turing reducibility conceptually;
    \item recognize the arithmetic hierarchy conceptually;
    \item understand computable approximations conceptually;
    \item connect computability with proof theory, recursion theory,
          complexity theory, type theory, and mathematical foundations.
\end{itemize}
}

%%%%%%%%%%%%%%%%%%%%%%%%%%%%%%%%%%%%%%%%%%%%%%%%%%%%%%%%%%%%
\subsection{What Is Computability Theory?}
\label{subsec:what-is-computability-theory}
%%%%%%%%%%%%%%%%%%%%%%%%%%%%%%%%%%%%%%%%%%%%%%%%%%%%%%%%%%%%

Computability theory studies which mathematical problems can be solved by
algorithms.

The central question is

\[
\boxed{
\text{What can be computed in principle?}
}
\]

This is different from asking

\[
\boxed{
\text{What can be computed efficiently?}
}
\]

The first question belongs primarily to computability theory.

The second belongs primarily to computational complexity.

%%%%%%%%%%%%%%%%%%%%%%%%%%%%%%%%%%%%%%%%%%%%%%%%%%%%%%%%%%%%
\subsection{Decision Problems}
\label{subsec:computability-decision-problems}
%%%%%%%%%%%%%%%%%%%%%%%%%%%%%%%%%%%%%%%%%%%%%%%%%%%%%%%%%%%%

A decision problem asks for a yes-or-no answer.

Examples include:

\[
\boxed{
\text{Is }n\text{ prime?}
}
\]

\[
\boxed{
\text{Does a graph contain a cycle?}
}
\]

and

\[
\boxed{
\text{Will a given program eventually halt?}
}
\]

Computability theory asks whether an algorithm exists that always produces
the correct answer.

%%%%%%%%%%%%%%%%%%%%%%%%%%%%%%%%%%%%%%%%%%%%%%%%%%%%%%%%%%%%
\subsection{Languages as Decision Problems}
\label{subsec:languages-as-decision-problems}
%%%%%%%%%%%%%%%%%%%%%%%%%%%%%%%%%%%%%%%%%%%%%%%%%%%%%%%%%%%%

Let

\[
\Sigma
\]

be a finite alphabet.

A language is a set

\[
\boxed{
L\subseteq\Sigma^*,
}
\]

where

\[
\Sigma^*
\]

is the set of all finite strings over \(\Sigma\).

A decision problem can be represented by asking:

\[
\boxed{
\text{Given }w\in\Sigma^*,
\text{ is }w\in L?
}
\]

Thus algorithms can be viewed as machines deciding languages.

%%%%%%%%%%%%%%%%%%%%%%%%%%%%%%%%%%%%%%%%%%%%%%%%%%%%%%%%%%%%
\subsection{Algorithms and Effective Procedures}
\label{subsec:effective-procedures}
%%%%%%%%%%%%%%%%%%%%%%%%%%%%%%%%%%%%%%%%%%%%%%%%%%%%%%%%%%%%

Informally, an algorithm is a finite, precise procedure for transforming
inputs into outputs.

A mathematical theory of computation requires a formal definition.

Several different formal models were developed, including:

\[
\boxed{
\text{Turing machines},
}
\]

\[
\boxed{
\lambda calculus},
\]

\[
\boxed{
\text{recursive functions},
}
\]

and

\[
\boxed{
\text{register machines}.
}
\]

These models agree on the broad class of effectively computable functions.

%%%%%%%%%%%%%%%%%%%%%%%%%%%%%%%%%%%%%%%%%%%%%%%%%%%%%%%%%%%%
\subsection{Turing Machines}
\label{subsec:turing-machines}
%%%%%%%%%%%%%%%%%%%%%%%%%%%%%%%%%%%%%%%%%%%%%%%%%%%%%%%%%%%%

A Turing machine is an idealized mathematical model of computation.

It consists conceptually of:

\[
\boxed{
\text{a finite set of states},
}
\]

\[
\boxed{
\text{an infinite tape},
}
\]

\[
\boxed{
\text{a read/write head},
}
\]

\[
\boxed{
\text{a finite tape alphabet},
}
\]

and

\[
\boxed{
\text{transition rules}.
}
\]

Despite its simplicity, the model captures general algorithmic computation.

%%%%%%%%%%%%%%%%%%%%%%%%%%%%%%%%%%%%%%%%%%%%%%%%%%%%%%%%%%%%
\subsection{Tape}
\label{subsec:turing-machine-tape}
%%%%%%%%%%%%%%%%%%%%%%%%%%%%%%%%%%%%%%%%%%%%%%%%%%%%%%%%%%%%

The tape is divided into cells.

Each cell stores one symbol from a finite alphabet

\[
\Gamma.
\]

One distinguished symbol is the blank symbol, often denoted

\[
\boxed{
\sqcup.
}
\]

Only finitely many tape cells are nonblank at the start of an ordinary
finite-input computation.

%%%%%%%%%%%%%%%%%%%%%%%%%%%%%%%%%%%%%%%%%%%%%%%%%%%%%%%%%%%%
\subsection{Read-Write Head}
\label{subsec:turing-machine-head}
%%%%%%%%%%%%%%%%%%%%%%%%%%%%%%%%%%%%%%%%%%%%%%%%%%%%%%%%%%%%

The machine has a head positioned over one tape cell.

At each step, the head:

\begin{enumerate}
    \item reads the current symbol;
    \item writes a symbol;
    \item changes state;
    \item moves one cell left or right, or possibly remains stationary
          depending on the chosen formalism.
\end{enumerate}

Thus local actions produce an arbitrarily long computation.

%%%%%%%%%%%%%%%%%%%%%%%%%%%%%%%%%%%%%%%%%%%%%%%%%%%%%%%%%%%%
\subsection{States}
\label{subsec:turing-machine-states}
%%%%%%%%%%%%%%%%%%%%%%%%%%%%%%%%%%%%%%%%%%%%%%%%%%%%%%%%%%%%

A Turing machine has a finite state set

\[
\boxed{
Q.
}
\]

Among these are typically:

\[
\boxed{
q_0
=
\text{start state},
}
\]

\[
\boxed{
q_{\mathrm{acc}}
=
\text{accept state},
}
\]

and

\[
\boxed{
q_{\mathrm{rej}}
=
\text{reject state}.
}
\]

A machine halts when it reaches a designated halting state.

%%%%%%%%%%%%%%%%%%%%%%%%%%%%%%%%%%%%%%%%%%%%%%%%%%%%%%%%%%%%
\subsection{Transition Function}
\label{subsec:turing-transition-function}
%%%%%%%%%%%%%%%%%%%%%%%%%%%%%%%%%%%%%%%%%%%%%%%%%%%%%%%%%%%%

For a deterministic Turing machine, the transition function has form

\[
\boxed{
\delta:
Q\times\Gamma
\longrightarrow
Q\times\Gamma\times\{L,R\}
}
\]

for nonhalting states.

If

\[
\delta(q,a)
=
(q',b,R),
\]

then the machine:

\begin{enumerate}
    \item currently in state \(q\);
    \item reads symbol \(a\);
    \item writes \(b\);
    \item changes to state \(q'\);
    \item moves one cell to the right.
\end{enumerate}

%%%%%%%%%%%%%%%%%%%%%%%%%%%%%%%%%%%%%%%%%%%%%%%%%%%%%%%%%%%%
\subsection{Formal Turing-Machine Description}
\label{subsec:formal-turing-machine}
%%%%%%%%%%%%%%%%%%%%%%%%%%%%%%%%%%%%%%%%%%%%%%%%%%%%%%%%%%%%

One common formalization writes a Turing machine as

\[
\boxed{
M
=
(Q,\Sigma,\Gamma,\delta,q_0,q_{\mathrm{acc}},q_{\mathrm{rej}}),
}
\]

where:

\[
\Sigma
\]

is the input alphabet,

\[
\Gamma
\]

is the tape alphabet,

and the remaining components specify the states and transitions.

Different textbooks use slightly different conventions without changing
the computational power.

%%%%%%%%%%%%%%%%%%%%%%%%%%%%%%%%%%%%%%%%%%%%%%%%%%%%%%%%%%%%
\subsection{Configurations}
\label{subsec:turing-configurations}
%%%%%%%%%%%%%%%%%%%%%%%%%%%%%%%%%%%%%%%%%%%%%%%%%%%%%%%%%%%%

A configuration records the complete instantaneous state of a machine.

It includes:

\[
\boxed{
\text{current state},
}
\]

\[
\boxed{
\text{tape contents},
}
\]

and

\[
\boxed{
\text{head position}.
}
\]

Thus a computation is a sequence

\[
\boxed{
C_0
\rightarrow
C_1
\rightarrow
C_2
\rightarrow
\cdots.
}
\]

%%%%%%%%%%%%%%%%%%%%%%%%%%%%%%%%%%%%%%%%%%%%%%%%%%%%%%%%%%%%
\subsection{Halting}
\label{subsec:turing-halting}
%%%%%%%%%%%%%%%%%%%%%%%%%%%%%%%%%%%%%%%%%%%%%%%%%%%%%%%%%%%%

A computation halts if it eventually reaches a designated halting state.

It may:

\[
\boxed{
\text{accept},
}
\]

\[
\boxed{
\text{reject},
}
\]

or, in some conventions, halt with a computed output.

A machine may also run forever.

Nontermination is a fundamental part of computability theory.

%%%%%%%%%%%%%%%%%%%%%%%%%%%%%%%%%%%%%%%%%%%%%%%%%%%%%%%%%%%%
\subsection{Recognizers}
\label{subsec:turing-recognizers}
%%%%%%%%%%%%%%%%%%%%%%%%%%%%%%%%%%%%%%%%%%%%%%%%%%%%%%%%%%%%

A Turing machine \(M\) recognizes language \(L\) if:

\[
w\in L
\Longrightarrow
M
\text{ eventually accepts }w,
\]

and

\[
w\notin L
\Longrightarrow
M
\text{ either rejects or runs forever}.
\]

Thus recognition guarantees halting only on positive instances.

%%%%%%%%%%%%%%%%%%%%%%%%%%%%%%%%%%%%%%%%%%%%%%%%%%%%%%%%%%%%
\subsection{Deciders}
\label{subsec:turing-deciders}
%%%%%%%%%%%%%%%%%%%%%%%%%%%%%%%%%%%%%%%%%%%%%%%%%%%%%%%%%%%%

A Turing machine decides \(L\) if it halts on every input and satisfies:

\[
\boxed{
w\in L
\Longrightarrow
M
\text{ accepts},
}
\]

and

\[
\boxed{
w\notin L
\Longrightarrow
M
\text{ rejects}.
}
\]

A decidable language is therefore algorithmically solvable.

%%%%%%%%%%%%%%%%%%%%%%%%%%%%%%%%%%%%%%%%%%%%%%%%%%%%%%%%%%%%
\subsection{Recognizable and Decidable Languages}
\label{subsec:recognizable-decidable}
%%%%%%%%%%%%%%%%%%%%%%%%%%%%%%%%%%%%%%%%%%%%%%%%%%%%%%%%%%%%

Every decidable language is recognizable:

\[
\boxed{
\text{decidable}
\Longrightarrow
\text{recognizable}.
}
\]

The converse is false.

There exist languages for which membership can be confirmed when true but
for which no algorithm always halts and answers correctly on every input.

%%%%%%%%%%%%%%%%%%%%%%%%%%%%%%%%%%%%%%%%%%%%%%%%%%%%%%%%%%%%
\subsection{Terminology}
\label{subsec:computability-terminology}
%%%%%%%%%%%%%%%%%%%%%%%%%%%%%%%%%%%%%%%%%%%%%%%%%%%%%%%%%%%%

Several equivalent terms are common.

A decidable set may be called:

\[
\boxed{
\text{recursive},
}
\]

or

\[
\boxed{
\text{computable}.
}
\]

A recognizable set may be called:

\[
\boxed{
\text{recursively enumerable},
}
\]

\[
\boxed{
\text{computably enumerable},
}
\]

or

\[
\boxed{
\text{semidecidable}.
}
\]

Modern usage often favors \textbf{computably enumerable}.

%%%%%%%%%%%%%%%%%%%%%%%%%%%%%%%%%%%%%%%%%%%%%%%%%%%%%%%%%%%%
\subsection{Worked Example: A Decidable Language}
\label{subsec:worked-decidable-language}
%%%%%%%%%%%%%%%%%%%%%%%%%%%%%%%%%%%%%%%%%%%%%%%%%%%%%%%%%%%%

\begin{workedexamplebox}{Even Binary Integers}

Let

\[
L
=
\{
w\in\{0,1\}^*:
w
\text{ represents an even binary integer}
\}.
\]

A binary integer is even exactly when its final digit is \(0\).

An algorithm therefore:

\begin{enumerate}
    \item reads the final input symbol;
    \item accepts if the symbol is \(0\);
    \item rejects otherwise.
\end{enumerate}

The procedure always halts.

\begin{answerbox}
\[
\boxed{
L
\text{ is decidable}.
}
\]
\end{answerbox}

\end{workedexamplebox}

%%%%%%%%%%%%%%%%%%%%%%%%%%%%%%%%%%%%%%%%%%%%%%%%%%%%%%%%%%%%
\subsection{Computable Functions}
\label{subsec:computable-functions}
%%%%%%%%%%%%%%%%%%%%%%%%%%%%%%%%%%%%%%%%%%%%%%%%%%%%%%%%%%%%

A function

\[
f:\Sigma^*\to\Sigma^*
\]

is computable if some Turing machine, on every input \(w\), halts with

\[
\boxed{
f(w)
}
\]

encoded on its tape.

Functions on natural numbers can be treated similarly through an encoding
of integers as strings.

%%%%%%%%%%%%%%%%%%%%%%%%%%%%%%%%%%%%%%%%%%%%%%%%%%%%%%%%%%%%
\subsection{Partial Computable Functions}
\label{subsec:partial-computable-functions}
%%%%%%%%%%%%%%%%%%%%%%%%%%%%%%%%%%%%%%%%%%%%%%%%%%%%%%%%%%%%

A partial function

\[
f:\Sigma^*
\rightharpoonup
\Sigma^*
\]

need not be defined on every input.

A machine computes \(f\) if:

\[
w\in\operatorname{dom}(f)
\Longrightarrow
M(w)
\text{ halts with output }f(w),
\]

while

\[
w\notin\operatorname{dom}(f)
\]

may cause the machine to run forever.

%%%%%%%%%%%%%%%%%%%%%%%%%%%%%%%%%%%%%%%%%%%%%%%%%%%%%%%%%%%%
\subsection{Total Versus Partial Computability}
\label{subsec:total-partial-computability}
%%%%%%%%%%%%%%%%%%%%%%%%%%%%%%%%%%%%%%%%%%%%%%%%%%%%%%%%%%%%

A total computable function halts on every valid input.

A partial computable function may fail to halt for some inputs.

Thus:

\[
\boxed{
\text{total computable}
\subsetneq
\text{partial computable}.
}
\]

Partial functions are essential because nontermination cannot be ignored in
a general theory of computation.

%%%%%%%%%%%%%%%%%%%%%%%%%%%%%%%%%%%%%%%%%%%%%%%%%%%%%%%%%%%%
\subsection{Characteristic Functions}
\label{subsec:characteristic-functions-computability}
%%%%%%%%%%%%%%%%%%%%%%%%%%%%%%%%%%%%%%%%%%%%%%%%%%%%%%%%%%%%

For a set

\[
A\subseteq\mathbb{N},
\]

define the characteristic function

\[
\boxed{
\chi_A(n)
=
\begin{cases}
1,& n\in A,\\
0,& n\notin A.
\end{cases}
}
\]

Then \(A\) is decidable exactly when

\[
\boxed{
\chi_A
}
\]

is a total computable function.

%%%%%%%%%%%%%%%%%%%%%%%%%%%%%%%%%%%%%%%%%%%%%%%%%%%%%%%%%%%%
\subsection{Enumerators}
\label{subsec:computability-enumerators}
%%%%%%%%%%%%%%%%%%%%%%%%%%%%%%%%%%%%%%%%%%%%%%%%%%%%%%%%%%%%

An enumerator is a machine that prints elements of a language, possibly
with repetitions and in arbitrary order.

A language \(L\) is computably enumerable if there exists an enumerator
whose set of outputs is exactly

\[
\boxed{
L.
}
\]

This provides another characterization of recognizability.

%%%%%%%%%%%%%%%%%%%%%%%%%%%%%%%%%%%%%%%%%%%%%%%%%%%%%%%%%%%%
\subsection{Recognizability and Enumeration}
\label{subsec:recognizability-enumeration}
%%%%%%%%%%%%%%%%%%%%%%%%%%%%%%%%%%%%%%%%%%%%%%%%%%%%%%%%%%%%

A language is Turing recognizable if and only if it can be enumerated by a
Turing machine.

Thus

\[
\boxed{
\text{recognizable}
\iff
\text{computably enumerable}.
}
\]

The proof uses systematic simulation of computations.

%%%%%%%%%%%%%%%%%%%%%%%%%%%%%%%%%%%%%%%%%%%%%%%%%%%%%%%%%%%%
\subsection{Dovetailing}
\label{subsec:dovetailing}
%%%%%%%%%%%%%%%%%%%%%%%%%%%%%%%%%%%%%%%%%%%%%%%%%%%%%%%%%%%%

Suppose infinitely many computations must be simulated, but some may never
halt.

Running them one at a time could get stuck forever.

Dovetailing interleaves the simulations:

\[
\boxed{
1\text{ step of computation }1,
}
\]

\[
\boxed{
1\text{ step each of computations }1,2,
}
\]

\[
\boxed{
1\text{ step each of computations }1,2,3,
}
\]

and so on.

No single nonhalting computation permanently blocks the others.

%%%%%%%%%%%%%%%%%%%%%%%%%%%%%%%%%%%%%%%%%%%%%%%%%%%%%%%%%%%%
\subsection{Worked Example: Dovetailing}
\label{subsec:worked-dovetailing}
%%%%%%%%%%%%%%%%%%%%%%%%%%%%%%%%%%%%%%%%%%%%%%%%%%%%%%%%%%%%

\begin{workedexamplebox}{Searching Many Computations}

Suppose machines

\[
M_1,M_2,M_3,\ldots
\]

are running on inputs

\[
w_1,w_2,w_3,\ldots.
\]

Instead of waiting for \(M_1(w_1)\) to finish, simulate:

\[
M_1
\]

for one step,

then

\[
M_1,M_2
\]

for another step each,

then

\[
M_1,M_2,M_3,
\]

and continue.

\begin{answerbox}
\[
\boxed{
\text{Every individual computation eventually receives arbitrarily many
simulation steps.}
}
\]
\end{answerbox}

\end{workedexamplebox}

%%%%%%%%%%%%%%%%%%%%%%%%%%%%%%%%%%%%%%%%%%%%%%%%%%%%%%%%%%%%
\subsection{Equivalent Models of Computation}
\label{subsec:equivalent-computation-models}
%%%%%%%%%%%%%%%%%%%%%%%%%%%%%%%%%%%%%%%%%%%%%%%%%%%%%%%%%%%%

Many mathematically different computational systems define the same class
of computable functions.

Examples include:

\[
\boxed{
\text{Turing machines},
}
\]

\[
\boxed{
\lambda calculus,
}
\]

\[
\boxed{
\mu\text{-recursive functions},
}
\]

and

\[
\boxed{
\text{register machines}.
}
\]

This robustness motivates the Church--Turing thesis.

%%%%%%%%%%%%%%%%%%%%%%%%%%%%%%%%%%%%%%%%%%%%%%%%%%%%%%%%%%%%
\subsection{Church--Turing Thesis}
\label{subsec:church-turing-thesis}
%%%%%%%%%%%%%%%%%%%%%%%%%%%%%%%%%%%%%%%%%%%%%%%%%%%%%%%%%%%%

The Church--Turing thesis states, informally:

\[
\boxed{
\text{every effectively calculable function is Turing computable}.
}
\]

This is not a theorem in the ordinary mathematical sense because
``effectively calculable'' is an informal concept.

It is a foundational thesis connecting intuitive algorithms with formal
computation.

%%%%%%%%%%%%%%%%%%%%%%%%%%%%%%%%%%%%%%%%%%%%%%%%%%%%%%%%%%%%
\subsection{Why the Church--Turing Thesis Is Plausible}
\label{subsec:why-church-turing}
%%%%%%%%%%%%%%%%%%%%%%%%%%%%%%%%%%%%%%%%%%%%%%%%%%%%%%%%%%%%

Its plausibility is supported by the fact that many independently proposed
models of effective computation produce the same class of functions.

Furthermore, ordinary programming languages can simulate Turing machines
and Turing machines can simulate general-purpose programs, abstracting away
finite resource limits.

%%%%%%%%%%%%%%%%%%%%%%%%%%%%%%%%%%%%%%%%%%%%%%%%%%%%%%%%%%%%
\subsection{Nondeterministic Turing Machines}
\label{subsec:nondeterministic-turing-machines}
%%%%%%%%%%%%%%%%%%%%%%%%%%%%%%%%%%%%%%%%%%%%%%%%%%%%%%%%%%%%

A nondeterministic Turing machine may have several possible transitions
from one configuration.

It accepts if at least one branch of computation accepts.

In computability theory, nondeterminism does not increase which languages
are decidable or recognizable:

\[
\boxed{
\text{deterministic and nondeterministic Turing machines have the same
computability power}.
}
\]

The difference becomes important in complexity theory.

%%%%%%%%%%%%%%%%%%%%%%%%%%%%%%%%%%%%%%%%%%%%%%%%%%%%%%%%%%%%
\subsection{Multitape Turing Machines}
\label{subsec:multitape-turing-machines}
%%%%%%%%%%%%%%%%%%%%%%%%%%%%%%%%%%%%%%%%%%%%%%%%%%%%%%%%%%%%

A multitape Turing machine has several tapes and heads.

It may be much more convenient or faster for some computations.

However, every multitape machine can be simulated by a standard single-tape
Turing machine.

Thus multitape machines do not change the class of computable functions.

%%%%%%%%%%%%%%%%%%%%%%%%%%%%%%%%%%%%%%%%%%%%%%%%%%%%%%%%%%%%
\subsection{Robustness of Computability}
\label{subsec:robustness-computability}
%%%%%%%%%%%%%%%%%%%%%%%%%%%%%%%%%%%%%%%%%%%%%%%%%%%%%%%%%%%%

Changing details such as:

\[
\boxed{
\text{number of tapes},
}
\]

\[
\boxed{
\text{alphabet size},
}
\]

or

\[
\boxed{
\text{determinism versus nondeterminism}
}
\]

does not fundamentally change which problems are computable.

This invariance is one of the strongest features of computability theory.

%%%%%%%%%%%%%%%%%%%%%%%%%%%%%%%%%%%%%%%%%%%%%%%%%%%%%%%%%%%%
\subsection{Encoding Machines}
\label{subsec:encoding-machines}
%%%%%%%%%%%%%%%%%%%%%%%%%%%%%%%%%%%%%%%%%%%%%%%%%%%%%%%%%%%%

A Turing machine has a finite description.

Therefore it can be encoded as a finite string or natural number.

Write

\[
\boxed{
\langle M\rangle
}
\]

for an encoding of machine \(M\).

Likewise,

\[
\boxed{
\langle M,w\rangle
}
\]

encodes a machine together with an input.

%%%%%%%%%%%%%%%%%%%%%%%%%%%%%%%%%%%%%%%%%%%%%%%%%%%%%%%%%%%%
\subsection{Countably Many Programs}
\label{subsec:countably-many-programs}
%%%%%%%%%%%%%%%%%%%%%%%%%%%%%%%%%%%%%%%%%%%%%%%%%%%%%%%%%%%%

Every Turing machine has a finite description over a finite alphabet.

Therefore the set of all Turing machines is countable.

One may list them:

\[
\boxed{
M_0,M_1,M_2,\ldots.
}
\]

The same holds for programs in any ordinary finite-description programming
language.

%%%%%%%%%%%%%%%%%%%%%%%%%%%%%%%%%%%%%%%%%%%%%%%%%%%%%%%%%%%%
\subsection{Uncountably Many Functions}
\label{subsec:uncountably-many-functions}
%%%%%%%%%%%%%%%%%%%%%%%%%%%%%%%%%%%%%%%%%%%%%%%%%%%%%%%%%%%%

Consider Boolean-valued functions

\[
f:\mathbb{N}\to\{0,1\}.
\]

Each such function corresponds to a subset of \(\mathbb{N}\).

Therefore there are

\[
\boxed{
2^{\aleph_0}
}
\]

such functions.

Since there are only countably many programs,

\[
\boxed{
\text{most functions }\mathbb{N}\to\{0,1\}
\text{ are not computable}.
}
\]

%%%%%%%%%%%%%%%%%%%%%%%%%%%%%%%%%%%%%%%%%%%%%%%%%%%%%%%%%%%%
\subsection{Universal Turing Machine}
\label{subsec:universal-turing-machine}
%%%%%%%%%%%%%%%%%%%%%%%%%%%%%%%%%%%%%%%%%%%%%%%%%%%%%%%%%%%%

A universal Turing machine \(U\) takes as input

\[
\boxed{
\langle M,w\rangle
}
\]

and simulates machine \(M\) on input \(w\).

Thus one machine can execute the descriptions of all other machines.

Conceptually,

\[
\boxed{
U(\langle M,w\rangle)
=
M(w).
}
\]

This is the theoretical basis of stored-program general-purpose computing.

%%%%%%%%%%%%%%%%%%%%%%%%%%%%%%%%%%%%%%%%%%%%%%%%%%%%%%%%%%%%
\subsection{Universal Computable Function}
\label{subsec:universal-computable-function}
%%%%%%%%%%%%%%%%%%%%%%%%%%%%%%%%%%%%%%%%%%%%%%%%%%%%%%%%%%%%

After enumerating programs

\[
M_0,M_1,\ldots,
\]

one can define a partial universal function

\[
\boxed{
U(e,x)
=
\varphi_e(x),
}
\]

where

\[
\varphi_e
\]

is the partial computable function computed by machine with index \(e\).

Universal computation allows programs to manipulate or simulate other
programs.

%%%%%%%%%%%%%%%%%%%%%%%%%%%%%%%%%%%%%%%%%%%%%%%%%%%%%%%%%%%%
\subsection{The Acceptance Problem}
\label{subsec:acceptance-problem}
%%%%%%%%%%%%%%%%%%%%%%%%%%%%%%%%%%%%%%%%%%%%%%%%%%%%%%%%%%%%

Define

\[
\boxed{
A_{\mathrm{TM}}
=
\{
\langle M,w\rangle:
M
\text{ accepts }w
\}.
}
\]

A universal simulator recognizes this language.

Given

\[
\langle M,w\rangle,
\]

simulate \(M\) on \(w\).

If \(M\) accepts, accept.

If \(M\) never halts, the simulation also never halts.

Thus

\[
\boxed{
A_{\mathrm{TM}}
\text{ is recognizable}.
}
\]

%%%%%%%%%%%%%%%%%%%%%%%%%%%%%%%%%%%%%%%%%%%%%%%%%%%%%%%%%%%%
\subsection{The Halting Problem}
\label{subsec:halting-problem}
%%%%%%%%%%%%%%%%%%%%%%%%%%%%%%%%%%%%%%%%%%%%%%%%%%%%%%%%%%%%

Define

\[
\boxed{
\operatorname{HALT}_{\mathrm{TM}}
=
\{
\langle M,w\rangle:
M
\text{ eventually halts on }w
\}.
}
\]

The Halting Problem asks whether there exists an algorithm deciding this
language.

The answer is

\[
\boxed{
\text{no}.
}
\]

%%%%%%%%%%%%%%%%%%%%%%%%%%%%%%%%%%%%%%%%%%%%%%%%%%%%%%%%%%%%
\subsection{Halting Is Recognizable}
\label{subsec:halting-recognizable}
%%%%%%%%%%%%%%%%%%%%%%%%%%%%%%%%%%%%%%%%%%%%%%%%%%%%%%%%%%%%

To recognize

\[
\operatorname{HALT}_{\mathrm{TM}},
\]

simulate \(M\) on \(w\).

If the simulation halts, accept.

If it runs forever, continue forever.

Therefore

\[
\boxed{
\operatorname{HALT}_{\mathrm{TM}}
\text{ is recognizable}.
}
\]

It is not decidable.

%%%%%%%%%%%%%%%%%%%%%%%%%%%%%%%%%%%%%%%%%%%%%%%%%%%%%%%%%%%%
\subsection{Proof of Halting Undecidability}
\label{subsec:halting-undecidability-proof}
%%%%%%%%%%%%%%%%%%%%%%%%%%%%%%%%%%%%%%%%%%%%%%%%%%%%%%%%%%%%

Assume for contradiction there exists a decider

\[
H(M,w)
\]

such that

\[
H(M,w)
=
\begin{cases}
\text{YES},&M\text{ halts on }w,\\
\text{NO},&M\text{ does not halt on }w.
\end{cases}
\]

Construct a new program \(D\) that takes a machine description
\(\langle M\rangle\) and does the opposite of the predicted behavior on
self-input:

\[
\boxed{
D(M)
=
\begin{cases}
\text{loop forever},&H(M,M)=\text{YES},\\
\text{halt},&H(M,M)=\text{NO}.
\end{cases}
}
\]

Now run \(D\) on its own description.

If \(H(D,D)\) says \(D\) halts, then \(D\) loops.

If \(H(D,D)\) says \(D\) loops, then \(D\) halts.

Either case is contradictory.

Therefore \(H\) cannot exist.

%%%%%%%%%%%%%%%%%%%%%%%%%%%%%%%%%%%%%%%%%%%%%%%%%%%%%%%%%%%%
\subsection{Worked Example: Diagonal Contradiction}
\label{subsec:worked-halting-diagonal}
%%%%%%%%%%%%%%%%%%%%%%%%%%%%%%%%%%%%%%%%%%%%%%%%%%%%%%%%%%%%

\begin{workedexamplebox}{Self-Application}

Suppose \(D\) is constructed so that

\[
D(M)
\]

halts exactly when the hypothetical halting decider predicts that \(M(M)\)
does not halt.

Now substitute

\[
M=D.
\]

Then

\[
D(D)
\]

halts if and only if

\[
D(D)
\]

does not halt.

\begin{answerbox}
\[
\boxed{
D(D)
\text{ halts}
\iff
D(D)
\text{ does not halt},
}
\]
\end{answerbox}

which is impossible.

Therefore a general halting decider cannot exist.

\end{workedexamplebox}

%%%%%%%%%%%%%%%%%%%%%%%%%%%%%%%%%%%%%%%%%%%%%%%%%%%%%%%%%%%%
\subsection{Diagonalization}
\label{subsec:computability-diagonalization}
%%%%%%%%%%%%%%%%%%%%%%%%%%%%%%%%%%%%%%%%%%%%%%%%%%%%%%%%%%%%

The Halting Problem proof is a computational version of diagonalization.

The broad pattern is:

\begin{enumerate}
    \item assume all objects of interest are uniformly captured;
    \item construct a new object that differs from the \(n\)-th object at
          position \(n\);
    \item derive contradiction.
\end{enumerate}

This resembles Cantor's diagonal proof and Gödel-style self-reference.

%%%%%%%%%%%%%%%%%%%%%%%%%%%%%%%%%%%%%%%%%%%%%%%%%%%%%%%%%%%%
\subsection{Recognizable but Undecidable}
\label{subsec:recognizable-undecidable}
%%%%%%%%%%%%%%%%%%%%%%%%%%%%%%%%%%%%%%%%%%%%%%%%%%%%%%%%%%%%

The Halting Problem is a canonical example of a language that is:

\[
\boxed{
\text{recognizable}
}
\]

but

\[
\boxed{
\text{not decidable}.
}
\]

Thus recognition is strictly weaker than decision.

%%%%%%%%%%%%%%%%%%%%%%%%%%%%%%%%%%%%%%%%%%%%%%%%%%%%%%%%%%%%
\subsection{Complement Languages}
\label{subsec:complement-languages-computability}
%%%%%%%%%%%%%%%%%%%%%%%%%%%%%%%%%%%%%%%%%%%%%%%%%%%%%%%%%%%%

For language

\[
L\subseteq\Sigma^*,
\]

define its complement by

\[
\boxed{
\overline{L}
=
\Sigma^*\setminus L.
}
\]

If \(L\) is decidable, then

\[
\boxed{
\overline{L}
}
\]

is also decidable.

Simply swap accept and reject outputs of the decider.

%%%%%%%%%%%%%%%%%%%%%%%%%%%%%%%%%%%%%%%%%%%%%%%%%%%%%%%%%%%%
\subsection{Recognizable and Co-Recognizable}
\label{subsec:co-recognizable}
%%%%%%%%%%%%%%%%%%%%%%%%%%%%%%%%%%%%%%%%%%%%%%%%%%%%%%%%%%%%

A language \(L\) is co-recognizable if

\[
\boxed{
\overline{L}
}
\]

is recognizable.

A fundamental theorem states:

\[
\boxed{
L
\text{ is decidable}
\iff
L
\text{ and }\overline{L}
\text{ are both recognizable}.
}
\]

%%%%%%%%%%%%%%%%%%%%%%%%%%%%%%%%%%%%%%%%%%%%%%%%%%%%%%%%%%%%
\subsection{Proof of the Recognition Characterization}
\label{subsec:recognizable-complement-proof}
%%%%%%%%%%%%%%%%%%%%%%%%%%%%%%%%%%%%%%%%%%%%%%%%%%%%%%%%%%%%

Suppose recognizer \(M_1\) recognizes \(L\) and recognizer \(M_2\) recognizes
\(\overline{L}\).

On input \(w\), dovetail the two simulations:

\[
\boxed{
M_1(w)
}
\]

and

\[
\boxed{
M_2(w).
}
\]

Exactly one must eventually accept.

If \(M_1\) accepts, accept.

If \(M_2\) accepts, reject.

Thus the combined machine halts on every input and decides \(L\).

%%%%%%%%%%%%%%%%%%%%%%%%%%%%%%%%%%%%%%%%%%%%%%%%%%%%%%%%%%%%
\subsection{The Complement of the Halting Problem}
\label{subsec:complement-halting}
%%%%%%%%%%%%%%%%%%%%%%%%%%%%%%%%%%%%%%%%%%%%%%%%%%%%%%%%%%%%

Since

\[
\operatorname{HALT}_{\mathrm{TM}}
\]

is recognizable but undecidable, its complement cannot also be
recognizable.

Otherwise the previous theorem would make the Halting Problem decidable.

Therefore

\[
\boxed{
\overline{
\operatorname{HALT}_{\mathrm{TM}}
}
}
\]

is not recognizable.

%%%%%%%%%%%%%%%%%%%%%%%%%%%%%%%%%%%%%%%%%%%%%%%%%%%%%%%%%%%%
\subsection{Closure of Decidable Languages}
\label{subsec:closure-decidable-languages}
%%%%%%%%%%%%%%%%%%%%%%%%%%%%%%%%%%%%%%%%%%%%%%%%%%%%%%%%%%%%

Decidable languages are closed under many Boolean operations.

If \(A\) and \(B\) are decidable, then so are:

\[
\boxed{
A\cup B,
}
\]

\[
\boxed{
A\cap B,
}
\]

\[
\boxed{
\overline{A},
}
\]

and

\[
\boxed{
A\setminus B.
}
\]

These results follow by combining deciders.

%%%%%%%%%%%%%%%%%%%%%%%%%%%%%%%%%%%%%%%%%%%%%%%%%%%%%%%%%%%%
\subsection{Closure of Recognizable Languages}
\label{subsec:closure-recognizable-languages}
%%%%%%%%%%%%%%%%%%%%%%%%%%%%%%%%%%%%%%%%%%%%%%%%%%%%%%%%%%%%

Recognizable languages are closed under:

\[
\boxed{
\text{union},
}
\]

and

\[
\boxed{
\text{intersection}.
}
\]

Dovetailing may be required.

They are not closed under complement.

Otherwise every recognizable language would be decidable.

%%%%%%%%%%%%%%%%%%%%%%%%%%%%%%%%%%%%%%%%%%%%%%%%%%%%%%%%%%%%
\subsection{Reductions}
\label{subsec:computability-reductions}
%%%%%%%%%%%%%%%%%%%%%%%%%%%%%%%%%%%%%%%%%%%%%%%%%%%%%%%%%%%%

A reduction transforms instances of one problem into instances of another.

If problem \(A\) can be transformed effectively into problem \(B\), then a
solver for \(B\) could be used to solve \(A\).

Thus reductions transfer impossibility results.

%%%%%%%%%%%%%%%%%%%%%%%%%%%%%%%%%%%%%%%%%%%%%%%%%%%%%%%%%%%%
\subsection{Mapping Reduction}
\label{subsec:mapping-reduction}
%%%%%%%%%%%%%%%%%%%%%%%%%%%%%%%%%%%%%%%%%%%%%%%%%%%%%%%%%%%%

Language \(A\) is mapping reducible to language \(B\), written

\[
\boxed{
A
\leq_m
B,
}
\]

if there exists a total computable function

\[
f:\Sigma^*\to\Sigma^*
\]

such that

\[
\boxed{
w\in A
\iff
f(w)\in B.
}
\]

The function \(f\) transforms membership questions about \(A\) into
membership questions about \(B\).

%%%%%%%%%%%%%%%%%%%%%%%%%%%%%%%%%%%%%%%%%%%%%%%%%%%%%%%%%%%%
\subsection{Using Reductions for Undecidability}
\label{subsec:reductions-undecidability}
%%%%%%%%%%%%%%%%%%%%%%%%%%%%%%%%%%%%%%%%%%%%%%%%%%%%%%%%%%%%

Suppose

\[
A
\]

is known undecidable and

\[
A\leq_m B.
\]

If \(B\) were decidable, then:

\begin{enumerate}
    \item compute \(f(w)\);
    \item decide whether \(f(w)\in B\);
    \item use equivalence to decide whether \(w\in A\).
\end{enumerate}

This would decide \(A\), contradiction.

Therefore,

\[
\boxed{
B
\text{ is undecidable}.
}
\]

%%%%%%%%%%%%%%%%%%%%%%%%%%%%%%%%%%%%%%%%%%%%%%%%%%%%%%%%%%%%
\subsection{Reduction Direction}
\label{subsec:reduction-direction}
%%%%%%%%%%%%%%%%%%%%%%%%%%%%%%%%%%%%%%%%%%%%%%%%%%%%%%%%%%%%

Reduction direction is crucial.

To prove \(B\) hard using known-hard problem \(A\), prove

\[
\boxed{
A
\leq_m
B.
}
\]

The known hard problem goes on the left.

A reduction

\[
B\leq_m A
\]

does not generally establish that \(B\) is hard.

%%%%%%%%%%%%%%%%%%%%%%%%%%%%%%%%%%%%%%%%%%%%%%%%%%%%%%%%%%%%
\subsection{Worked Example: Reduction Logic}
\label{subsec:worked-reduction-logic}
%%%%%%%%%%%%%%%%%%%%%%%%%%%%%%%%%%%%%%%%%%%%%%%%%%%%%%%%%%%%

\begin{workedexamplebox}{Why the Direction Matters}

Suppose \(A\) is undecidable and

\[
A\leq_m B.
\]

Assume \(B\) has a decider.

Then transform input \(w\) for \(A\) into

\[
f(w)
\]

and run the \(B\)-decider.

Because

\[
w\in A
\iff
f(w)\in B,
\]

this decides \(A\).

Contradiction.

\begin{answerbox}
\[
\boxed{
A\leq_m B
\text{ and }
A
\text{ undecidable}
\Longrightarrow
B
\text{ undecidable}.
}
\]
\end{answerbox}

\end{workedexamplebox}

%%%%%%%%%%%%%%%%%%%%%%%%%%%%%%%%%%%%%%%%%%%%%%%%%%%%%%%%%%%%
\subsection{Acceptance Problem Undecidability}
\label{subsec:acceptance-undecidability}
%%%%%%%%%%%%%%%%%%%%%%%%%%%%%%%%%%%%%%%%%%%%%%%%%%%%%%%%%%%%

The language

\[
A_{\mathrm{TM}}
=
\{
\langle M,w\rangle:
M
\text{ accepts }w
\}
\]

is recognizable but undecidable.

One can prove this directly by diagonalization or by reduction from the
Halting Problem.

Thus even determining whether a program eventually accepts a particular
input is undecidable in general.

%%%%%%%%%%%%%%%%%%%%%%%%%%%%%%%%%%%%%%%%%%%%%%%%%%%%%%%%%%%%
\subsection{Emptiness Problem for Turing Machines}
\label{subsec:turing-emptiness}
%%%%%%%%%%%%%%%%%%%%%%%%%%%%%%%%%%%%%%%%%%%%%%%%%%%%%%%%%%%%

Define

\[
\boxed{
E_{\mathrm{TM}}
=
\{
\langle M\rangle:
L(M)=\varnothing
\}.
}
\]

This asks whether machine \(M\) accepts no strings.

The problem is undecidable.

The property concerns the entire language recognized by \(M\), not merely
one computation.

%%%%%%%%%%%%%%%%%%%%%%%%%%%%%%%%%%%%%%%%%%%%%%%%%%%%%%%%%%%%
\subsection{Equivalence Problem}
\label{subsec:turing-equivalence-problem}
%%%%%%%%%%%%%%%%%%%%%%%%%%%%%%%%%%%%%%%%%%%%%%%%%%%%%%%%%%%%

Define

\[
\boxed{
EQ_{\mathrm{TM}}
=
\{
\langle M_1,M_2\rangle:
L(M_1)=L(M_2)
\}.
}
\]

This asks whether two arbitrary Turing machines recognize exactly the same
language.

The problem is undecidable.

General program equivalence is therefore impossible to decide
algorithmically.

%%%%%%%%%%%%%%%%%%%%%%%%%%%%%%%%%%%%%%%%%%%%%%%%%%%%%%%%%%%%
\subsection{Rice's Theorem}
\label{subsec:rices-theorem}
%%%%%%%%%%%%%%%%%%%%%%%%%%%%%%%%%%%%%%%%%%%%%%%%%%%%%%%%%%%%

\begin{theorem}
Every nontrivial semantic property of the language recognized by a Turing
machine is undecidable.
\end{theorem}

A property is nontrivial if some recognizable languages possess it and some
do not.

Thus properties such as

\[
\boxed{
L(M)=\varnothing,
}
\]

\[
\boxed{
L(M)
\text{ is finite},
}
\]

or

\[
\boxed{
L(M)
\text{ contains a particular string}
}
\]

are generally undecidable when phrased as semantic properties of arbitrary
Turing-machine languages.

%%%%%%%%%%%%%%%%%%%%%%%%%%%%%%%%%%%%%%%%%%%%%%%%%%%%%%%%%%%%
\subsection{Syntactic Versus Semantic Program Properties}
\label{subsec:syntactic-semantic-program-properties}
%%%%%%%%%%%%%%%%%%%%%%%%%%%%%%%%%%%%%%%%%%%%%%%%%%%%%%%%%%%%

Rice's theorem applies to semantic behavior, not arbitrary syntactic
properties of program descriptions.

For example,

\[
\boxed{
\text{``Does this program description contain 100 states?''}
}
\]

is syntactic and decidable from the encoding.

By contrast,

\[
\boxed{
\text{``Does this program compute the zero function?''}
}
\]

is semantic and nontrivial, so it falls under Rice-style undecidability.

%%%%%%%%%%%%%%%%%%%%%%%%%%%%%%%%%%%%%%%%%%%%%%%%%%%%%%%%%%%%
\subsection{Worked Example: Rice's Theorem}
\label{subsec:worked-rice}
%%%%%%%%%%%%%%%%%%%%%%%%%%%%%%%%%%%%%%%%%%%%%%%%%%%%%%%%%%%%

\begin{workedexamplebox}{Does a Machine Accept Anything?}

Consider the property

\[
P(M):
L(M)\neq\varnothing.
\]

Some machines accept at least one string.

Others accept none.

Thus the property is nontrivial.

It depends only on the language recognized, not on the machine's syntax.

Therefore Rice's theorem implies:

\begin{answerbox}
\[
\boxed{
\text{The property }L(M)\neq\varnothing
\text{ is undecidable}.
}
\]
\end{answerbox}

\end{workedexamplebox}

%%%%%%%%%%%%%%%%%%%%%%%%%%%%%%%%%%%%%%%%%%%%%%%%%%%%%%%%%%%%
\subsection{Post Correspondence Problem}
\label{subsec:post-correspondence-problem}
%%%%%%%%%%%%%%%%%%%%%%%%%%%%%%%%%%%%%%%%%%%%%%%%%%%%%%%%%%%%

The Post Correspondence Problem, abbreviated PCP, is a classical
undecidable string-matching problem.

One is given pairs

\[
\boxed{
(u_1,v_1),
\ldots,
(u_n,v_n).
}
\]

The question is whether there exists a nonempty index sequence

\[
i_1,\ldots,i_k
\]

such that

\[
\boxed{
u_{i_1}\cdots u_{i_k}
=
v_{i_1}\cdots v_{i_k}.
}
\]

PCP provides a useful source problem for reductions involving formal
languages.

%%%%%%%%%%%%%%%%%%%%%%%%%%%%%%%%%%%%%%%%%%%%%%%%%%%%%%%%%%%%
\subsection{Worked Example: A PCP Match}
\label{subsec:worked-pcp}
%%%%%%%%%%%%%%%%%%%%%%%%%%%%%%%%%%%%%%%%%%%%%%%%%%%%%%%%%%%%

\begin{workedexamplebox}{Matching Two String Sequences}

Suppose the available pairs are

\[
(a,ab)
\]

and

\[
(ba,a).
\]

Choosing the sequence

\[
1,2
\]

gives top string

\[
a\,ba
=
aba,
\]

and bottom string

\[
ab\,a
=
aba.
\]

Therefore,

\begin{answerbox}
\[
\boxed{
1,2
\text{ is a PCP solution for this instance}.
}
\]
\end{answerbox}

\end{workedexamplebox}

%%%%%%%%%%%%%%%%%%%%%%%%%%%%%%%%%%%%%%%%%%%%%%%%%%%%%%%%%%%%
\subsection{Bounded Simulation}
\label{subsec:bounded-simulation}
%%%%%%%%%%%%%%%%%%%%%%%%%%%%%%%%%%%%%%%%%%%%%%%%%%%%%%%%%%%%

Although general halting is undecidable, bounded halting is decidable.

Given:

\[
\boxed{
M,
\quad
w,
\quad
t,
}
\]

simulate \(M(w)\) for at most \(t\) steps.

Then determine whether the machine halted within that bound.

Thus

\[
\boxed{
\text{unbounded future behavior}
}
\]

can be undecidable even when every finite bounded simulation is
straightforward.

%%%%%%%%%%%%%%%%%%%%%%%%%%%%%%%%%%%%%%%%%%%%%%%%%%%%%%%%%%%%
\subsection{Finite Versus Infinite Search}
\label{subsec:finite-infinite-search}
%%%%%%%%%%%%%%%%%%%%%%%%%%%%%%%%%%%%%%%%%%%%%%%%%%%%%%%%%%%%

A finite search can always be exhausted in principle.

For example, checking all strings of length at most \(n\) is finite.

An unbounded search may fail because there is no computable point at which
one can conclude that a desired witness will never appear.

This distinction underlies many recognizable but undecidable problems.

%%%%%%%%%%%%%%%%%%%%%%%%%%%%%%%%%%%%%%%%%%%%%%%%%%%%%%%%%%%%
\subsection{Witness Search}
\label{subsec:computable-witness-search}
%%%%%%%%%%%%%%%%%%%%%%%%%%%%%%%%%%%%%%%%%%%%%%%%%%%%%%%%%%%%

Suppose membership in \(L\) means:

\[
\boxed{
w\in L
\iff
\exists y\,R(w,y),
}
\]

where \(R\) is decidable.

One can systematically search through all possible \(y\).

If a witness exists, it will eventually be found.

If no witness exists, the search may continue forever.

Thus such sets are natural candidates for computable enumerability.

%%%%%%%%%%%%%%%%%%%%%%%%%%%%%%%%%%%%%%%%%%%%%%%%%%%%%%%%%%%%
\subsection{Effective Enumerations}
\label{subsec:effective-enumerations}
%%%%%%%%%%%%%%%%%%%%%%%%%%%%%%%%%%%%%%%%%%%%%%%%%%%%%%%%%%%%

A sequence

\[
a_0,a_1,a_2,\ldots
\]

is effectively generated if there is a computable function

\[
\boxed{
f(n)=a_n.
}
\]

Effective enumeration is stronger than mere countability.

A set can be countable without having an effective enumeration available
within a particular formal framework.

%%%%%%%%%%%%%%%%%%%%%%%%%%%%%%%%%%%%%%%%%%%%%%%%%%%%%%%%%%%%
\subsection{Pairing Functions}
\label{subsec:computable-pairing-functions}
%%%%%%%%%%%%%%%%%%%%%%%%%%%%%%%%%%%%%%%%%%%%%%%%%%%%%%%%%%%%

Computability theory often encodes tuples of natural numbers as single
natural numbers.

A pairing function is a computable bijection or effective coding

\[
\boxed{
\langle\cdot,\cdot\rangle:
\mathbb{N}^2
\rightarrow
\mathbb{N}.
}
\]

The projections should also be computable.

This lets multivariable problems be encoded as unary ones.

%%%%%%%%%%%%%%%%%%%%%%%%%%%%%%%%%%%%%%%%%%%%%%%%%%%%%%%%%%%%
\subsection{Coding Finite Sequences}
\label{subsec:coding-finite-sequences}
%%%%%%%%%%%%%%%%%%%%%%%%%%%%%%%%%%%%%%%%%%%%%%%%%%%%%%%%%%%%

Finite sequences

\[
(a_0,\ldots,a_n)
\]

can also be encoded by natural numbers.

Thus:

\[
\boxed{
\text{strings},
}
\]

\[
\boxed{
\text{machine descriptions},
}
\]

\[
\boxed{
\text{proofs},
}
\]

and

\[
\boxed{
\text{computation histories}
}
\]

can all be represented arithmetically.

This parallels Gödel coding in proof theory.

%%%%%%%%%%%%%%%%%%%%%%%%%%%%%%%%%%%%%%%%%%%%%%%%%%%%%%%%%%%%
\subsection{Computation Histories}
\label{subsec:computation-histories}
%%%%%%%%%%%%%%%%%%%%%%%%%%%%%%%%%%%%%%%%%%%%%%%%%%%%%%%%%%%%

A halting computation can be represented as a finite sequence of
configurations:

\[
\boxed{
C_0,C_1,\ldots,C_t.
}
\]

Each adjacent pair must satisfy the machine's transition relation.

Because the history is finite, a proposed halting computation can be
verified mechanically.

This is important in both computability and complexity theory.

%%%%%%%%%%%%%%%%%%%%%%%%%%%%%%%%%%%%%%%%%%%%%%%%%%%%%%%%%%%%
\subsection{Universal Search}
\label{subsec:universal-search-computability}
%%%%%%%%%%%%%%%%%%%%%%%%%%%%%%%%%%%%%%%%%%%%%%%%%%%%%%%%%%%%

Suppose a property has mechanically verifiable witnesses.

One can enumerate all possible witnesses and test them.

This gives a universal search procedure.

The search may be inefficient, but if a witness exists it is eventually
found.

Thus computability ignores efficiency unless efficiency is explicitly part
of the problem.

%%%%%%%%%%%%%%%%%%%%%%%%%%%%%%%%%%%%%%%%%%%%%%%%%%%%%%%%%%%%
\subsection{Computability Versus Complexity}
\label{subsec:computability-versus-complexity}
%%%%%%%%%%%%%%%%%%%%%%%%%%%%%%%%%%%%%%%%%%%%%%%%%%%%%%%%%%%%

Computability asks:

\[
\boxed{
\text{Does any terminating algorithm exist?}
}
\]

Complexity asks:

\[
\boxed{
\text{How many resources does an algorithm require?}
}
\]

A problem may be decidable but require enormous computational resources.

Therefore:

\[
\boxed{
\text{hard}
\neq
\text{undecidable}.
}
\]

%%%%%%%%%%%%%%%%%%%%%%%%%%%%%%%%%%%%%%%%%%%%%%%%%%%%%%%%%%%%
\subsection{Decidable but Expensive Problems}
\label{subsec:decidable-expensive-problems}
%%%%%%%%%%%%%%%%%%%%%%%%%%%%%%%%%%%%%%%%%%%%%%%%%%%%%%%%%%%%

A finite combinatorial problem may be solvable by exhaustive search.

For example, a problem over \(n\) Boolean variables can be decided by
checking all

\[
\boxed{
2^n
}
\]

assignments.

This proves decidability.

It does not imply practical efficiency.

%%%%%%%%%%%%%%%%%%%%%%%%%%%%%%%%%%%%%%%%%%%%%%%%%%%%%%%%%%%%
\subsection{Undecidability Is Stronger Than Intractability}
\label{subsec:undecidability-vs-intractability}
%%%%%%%%%%%%%%%%%%%%%%%%%%%%%%%%%%%%%%%%%%%%%%%%%%%%%%%%%%%%

For an undecidable problem, no algorithm can correctly halt on every input.

This is stronger than merely saying known algorithms are slow.

No amount of faster hardware eliminates the theoretical impossibility.

%%%%%%%%%%%%%%%%%%%%%%%%%%%%%%%%%%%%%%%%%%%%%%%%%%%%%%%%%%%%
\subsection{Oracle Computation}
\label{subsec:oracle-computation}
%%%%%%%%%%%%%%%%%%%%%%%%%%%%%%%%%%%%%%%%%%%%%%%%%%%%%%%%%%%%

An oracle Turing machine is an ordinary Turing machine augmented with
access to a hypothetical black box that instantly answers membership
questions for some set \(A\).

One writes conceptually

\[
\boxed{
M^A.
}
\]

The oracle changes which problems the machine can solve.

%%%%%%%%%%%%%%%%%%%%%%%%%%%%%%%%%%%%%%%%%%%%%%%%%%%%%%%%%%%%
\subsection{Relative Computability}
\label{subsec:relative-computability}
%%%%%%%%%%%%%%%%%%%%%%%%%%%%%%%%%%%%%%%%%%%%%%%%%%%%%%%%%%%%

A set \(B\) is computable relative to \(A\) if an oracle machine using \(A\)
can decide \(B\).

This does not mean \(B\) is computable without the oracle.

Relative computability compares the informational strength of different
problems.

%%%%%%%%%%%%%%%%%%%%%%%%%%%%%%%%%%%%%%%%%%%%%%%%%%%%%%%%%%%%
\subsection{Turing Reducibility}
\label{subsec:turing-reducibility}
%%%%%%%%%%%%%%%%%%%%%%%%%%%%%%%%%%%%%%%%%%%%%%%%%%%%%%%%%%%%

Write

\[
\boxed{
A
\leq_T
B
}
\]

if \(A\) can be decided by a Turing machine with oracle access to \(B\).

Turing reducibility is more flexible than mapping reducibility because the
machine may query the oracle repeatedly and adaptively.

%%%%%%%%%%%%%%%%%%%%%%%%%%%%%%%%%%%%%%%%%%%%%%%%%%%%%%%%%%%%
\subsection{Turing Degrees Preview}
\label{subsec:turing-degrees-preview}
%%%%%%%%%%%%%%%%%%%%%%%%%%%%%%%%%%%%%%%%%%%%%%%%%%%%%%%%%%%%

Two sets have the same Turing degree if each is Turing reducible to the
other:

\[
\boxed{
A\leq_TB
\quad\text{and}\quad
B\leq_TA.
}
\]

A Turing degree therefore represents a level of computational
unsolvability.

The computable sets form the lowest degree.

%%%%%%%%%%%%%%%%%%%%%%%%%%%%%%%%%%%%%%%%%%%%%%%%%%%%%%%%%%%%
\subsection{The Halting Degree}
\label{subsec:halting-degree}
%%%%%%%%%%%%%%%%%%%%%%%%%%%%%%%%%%%%%%%%%%%%%%%%%%%%%%%%%%%%

The Halting Problem is not computable.

However, if a machine has access to a halting oracle, it can solve many
problems that ordinary Turing machines cannot.

Thus the halting set occupies a strictly higher level of relative
computability than computable sets.

%%%%%%%%%%%%%%%%%%%%%%%%%%%%%%%%%%%%%%%%%%%%%%%%%%%%%%%%%%%%
\subsection{The Turing Jump}
\label{subsec:turing-jump-preview}
%%%%%%%%%%%%%%%%%%%%%%%%%%%%%%%%%%%%%%%%%%%%%%%%%%%%%%%%%%%%

Given oracle \(A\), one may define a new halting problem for machines with
oracle access to \(A\).

This produces the Turing jump

\[
\boxed{
A'.
}
\]

The jump is strictly more powerful in the sense that

\[
\boxed{
A
<_T
A'.
}
\]

Repeated jumps generate a hierarchy of increasing computational strength.

%%%%%%%%%%%%%%%%%%%%%%%%%%%%%%%%%%%%%%%%%%%%%%%%%%%%%%%%%%%%
\subsection{Arithmetic Hierarchy Preview}
\label{subsec:arithmetic-hierarchy-preview}
%%%%%%%%%%%%%%%%%%%%%%%%%%%%%%%%%%%%%%%%%%%%%%%%%%%%%%%%%%%%

The arithmetic hierarchy classifies subsets of \(\mathbb{N}\) by the
logical complexity of number quantifiers in their definitions.

The levels are often written:

\[
\boxed{
\Sigma^0_1,
\Pi^0_1,
\Sigma^0_2,
\Pi^0_2,
\ldots.
}
\]

Computably enumerable sets correspond closely to the first existential
level.

This topic is developed further in recursion theory.

%%%%%%%%%%%%%%%%%%%%%%%%%%%%%%%%%%%%%%%%%%%%%%%%%%%%%%%%%%%%
\subsection{\(\Sigma^0_1\) Form}
\label{subsec:sigma-one-form}
%%%%%%%%%%%%%%%%%%%%%%%%%%%%%%%%%%%%%%%%%%%%%%%%%%%%%%%%%%%%

A typical \(\Sigma^0_1\) property has form

\[
\boxed{
n\in A
\iff
\exists y\,R(n,y),
}
\]

where \(R\) is computable.

The witness \(y\) can be searched for effectively.

Therefore such sets are computably enumerable.

%%%%%%%%%%%%%%%%%%%%%%%%%%%%%%%%%%%%%%%%%%%%%%%%%%%%%%%%%%%%
\subsection{\(\Pi^0_1\) Form}
\label{subsec:pi-one-form}
%%%%%%%%%%%%%%%%%%%%%%%%%%%%%%%%%%%%%%%%%%%%%%%%%%%%%%%%%%%%

A typical \(\Pi^0_1\) property has form

\[
\boxed{
n\in A
\iff
\forall y\,R(n,y),
}
\]

where \(R\) is computable.

A counterexample can be found if one exists, but verifying that no
counterexample will ever appear may be impossible.

This asymmetry is central to semidecidability.

%%%%%%%%%%%%%%%%%%%%%%%%%%%%%%%%%%%%%%%%%%%%%%%%%%%%%%%%%%%%
\subsection{Program Verification Limits}
\label{subsec:program-verification-limits}
%%%%%%%%%%%%%%%%%%%%%%%%%%%%%%%%%%%%%%%%%%%%%%%%%%%%%%%%%%%%

Undecidability does not mean that no individual programs can be verified.

Many programs have provable properties.

Rather, there is no universal algorithm that correctly decides every
nontrivial semantic property of every arbitrary program.

Thus practical verification relies on restricted languages, annotations,
proofs, abstractions, or incomplete procedures.

%%%%%%%%%%%%%%%%%%%%%%%%%%%%%%%%%%%%%%%%%%%%%%%%%%%%%%%%%%%%
\subsection{Termination Analysis}
\label{subsec:termination-analysis}
%%%%%%%%%%%%%%%%%%%%%%%%%%%%%%%%%%%%%%%%%%%%%%%%%%%%%%%%%%%%

For a particular program, termination may be proved using a decreasing
measure.

For example, if each recursive call decreases

\[
n\in\mathbb{N}
\]

and cannot continue below zero, termination follows.

The Halting Problem says there is no algorithm that successfully decides
termination for every possible program-input pair.

%%%%%%%%%%%%%%%%%%%%%%%%%%%%%%%%%%%%%%%%%%%%%%%%%%%%%%%%%%%%
\subsection{Worked Example: Proving Termination}
\label{subsec:worked-termination-measure}
%%%%%%%%%%%%%%%%%%%%%%%%%%%%%%%%%%%%%%%%%%%%%%%%%%%%%%%%%%%%

\begin{workedexamplebox}{A Decreasing Natural Measure}

Suppose an algorithm repeatedly replaces positive integer \(n\) by

\[
n-1.
\]

The quantity

\[
n
\]

strictly decreases at each step and remains in

\[
\mathbb{N}.
\]

There is no infinite strictly descending sequence

\[
n>n-1>n-2>\cdots
\]

of natural numbers.

\begin{answerbox}
\[
\boxed{
\text{The algorithm must terminate after at most the initial }n
\text{ steps}.
}
\]
\end{answerbox}

\end{workedexamplebox}

%%%%%%%%%%%%%%%%%%%%%%%%%%%%%%%%%%%%%%%%%%%%%%%%%%%%%%%%%%%%
\subsection{Computability and Formal Proof}
\label{subsec:computability-formal-proof}
%%%%%%%%%%%%%%%%%%%%%%%%%%%%%%%%%%%%%%%%%%%%%%%%%%%%%%%%%%%%

Formal proofs can be mechanically enumerated when the proof system is
effectively specified.

Thus theorems of an effectively axiomatized formal theory are typically
computably enumerable.

One may systematically generate candidate finite strings and test which
ones are valid proofs.

%%%%%%%%%%%%%%%%%%%%%%%%%%%%%%%%%%%%%%%%%%%%%%%%%%%%%%%%%%%%
\subsection{Theoremhood and Recognizability}
\label{subsec:theoremhood-recognizability}
%%%%%%%%%%%%%%%%%%%%%%%%%%%%%%%%%%%%%%%%%%%%%%%%%%%%%%%%%%%%

Let

\[
\operatorname{Thm}(T)
=
\{
\varphi:
T\vdash\varphi
\}.
\]

If \(T\) is effectively axiomatized and proofs are mechanically checkable,
then

\[
\boxed{
\operatorname{Thm}(T)
}
\]

is computably enumerable.

Given a formula \(\varphi\), enumerate all \(T\)-proofs until one proves
\(\varphi\).

If no proof exists, the search may continue forever.

%%%%%%%%%%%%%%%%%%%%%%%%%%%%%%%%%%%%%%%%%%%%%%%%%%%%%%%%%%%%
\subsection{Decidable Theories}
\label{subsec:decidable-theories-computability}
%%%%%%%%%%%%%%%%%%%%%%%%%%%%%%%%%%%%%%%%%%%%%%%%%%%%%%%%%%%%

A theory \(T\) is decidable if there exists an algorithm that, for every
sentence \(\varphi\), determines whether

\[
\boxed{
T\models\varphi.
}
\]

Some important mathematical theories are decidable.

Others, particularly sufficiently expressive arithmetic theories, are
undecidable.

%%%%%%%%%%%%%%%%%%%%%%%%%%%%%%%%%%%%%%%%%%%%%%%%%%%%%%%%%%%%
\subsection{Undecidable Theories}
\label{subsec:undecidable-theories}
%%%%%%%%%%%%%%%%%%%%%%%%%%%%%%%%%%%%%%%%%%%%%%%%%%%%%%%%%%%%

A theory may be complete in the semantic sense yet still fail to possess an
effective decision algorithm.

Completeness and decidability are therefore distinct properties.

Similarly, effective axiomatizability and decidability are different.

%%%%%%%%%%%%%%%%%%%%%%%%%%%%%%%%%%%%%%%%%%%%%%%%%%%%%%%%%%%%
\subsection{Computability and Self-Reference}
\label{subsec:computability-self-reference}
%%%%%%%%%%%%%%%%%%%%%%%%%%%%%%%%%%%%%%%%%%%%%%%%%%%%%%%%%%%%

Programs can operate on encodings of programs.

This permits constructions in which a program is applied to its own
description.

Such self-reference is central to:

\[
\boxed{
\text{Halting Problem proofs},
}
\]

\[
\boxed{
\text{recursion theorems},
}
\]

and

\[
\boxed{
\text{Gödel-style incompleteness}.
}
\]

%%%%%%%%%%%%%%%%%%%%%%%%%%%%%%%%%%%%%%%%%%%%%%%%%%%%%%%%%%%%
\subsection{Fixed-Point Phenomena Preview}
\label{subsec:computability-fixed-point-preview}
%%%%%%%%%%%%%%%%%%%%%%%%%%%%%%%%%%%%%%%%%%%%%%%%%%%%%%%%%%%%

Computability theory contains formal fixed-point theorems showing that
programs can, in a precise sense, obtain or reproduce descriptions of their
own behavior.

These ideas lead to the recursion theorem developed more deeply in the
next section.

%%%%%%%%%%%%%%%%%%%%%%%%%%%%%%%%%%%%%%%%%%%%%%%%%%%%%%%%%%%%
\subsection{Computability Workflow}
\label{subsec:computability-workflow}
%%%%%%%%%%%%%%%%%%%%%%%%%%%%%%%%%%%%%%%%%%%%%%%%%%%%%%%%%%%%

When analyzing a computability problem:

\begin{enumerate}
    \item define the input encoding precisely;
    \item identify the output or yes-no question;
    \item determine whether the goal is recognition or decision;
    \item attempt to construct a terminating algorithm;
    \item separate finite search from unbounded search;
    \item identify whether witnesses can be effectively verified;
    \item use dovetailing when several nonterminating searches must be
          combined;
    \item determine whether the complement is recognizable;
    \item identify known undecidable problems that may reduce to the target;
    \item choose the correct reduction direction;
    \item construct the reduction algorithm explicitly;
    \item prove both directions of the membership equivalence;
    \item use diagonalization when self-reference is natural;
    \item use Rice's theorem for nontrivial semantic properties of arbitrary
          Turing-machine languages;
    \item distinguish syntactic properties from semantic properties;
    \item distinguish undecidability from high computational complexity;
    \item consider relative computability when oracles are introduced;
    \item state the computational model clearly;
    \item verify that all encodings and transformations are computable.
\end{enumerate}

%%%%%%%%%%%%%%%%%%%%%%%%%%%%%%%%%%%%%%%%%%%%%%%%%%%%%%%%%%%%
\subsection{Common Errors}
\label{subsec:computability-common-errors}
%%%%%%%%%%%%%%%%%%%%%%%%%%%%%%%%%%%%%%%%%%%%%%%%%%%%%%%%%%%%

\subsubsection{Confusing Recognizable With Decidable}

A recognizer may run forever on negative instances.

A decider must halt on every input.

Thus

\[
\boxed{
\text{recognizable}
\not\Rightarrow
\text{decidable}.
}
\]

%%%%%%%%%%%%%%%%%%%%%%%%%%%%%%%%%%%%%%%%%%%%%%%%%%%%%%%%%%%%
\subsubsection{Assuming Nontermination Means Rejection}

If a machine loops forever, it has not rejected unless the model explicitly
defines a rejecting halt.

Infinite computation and rejecting computation are different.

%%%%%%%%%%%%%%%%%%%%%%%%%%%%%%%%%%%%%%%%%%%%%%%%%%%%%%%%%%%%
\subsubsection{Confusing Undecidable With Unrecognizable}

Some undecidable problems are recognizable.

The Halting Problem is the standard example.

%%%%%%%%%%%%%%%%%%%%%%%%%%%%%%%%%%%%%%%%%%%%%%%%%%%%%%%%%%%%
\subsubsection{Assuming Every Countable Set Is Computably Enumerable}

Countability means an enumeration exists set-theoretically.

Computable enumerability requires an effective enumeration.

These are different notions.

%%%%%%%%%%%%%%%%%%%%%%%%%%%%%%%%%%%%%%%%%%%%%%%%%%%%%%%%%%%%
\subsubsection{Confusing Computability With Practical Feasibility}

A computation may terminate only after an astronomically large number of
steps.

It is still computable in the theoretical sense.

%%%%%%%%%%%%%%%%%%%%%%%%%%%%%%%%%%%%%%%%%%%%%%%%%%%%%%%%%%%%
\subsubsection{Thinking Turing Machines Are Claims About Physical Hardware}

A Turing machine is a mathematical model.

It abstracts the notion of effective computation.

Its purpose is foundational, not mechanical realism.

%%%%%%%%%%%%%%%%%%%%%%%%%%%%%%%%%%%%%%%%%%%%%%%%%%%%%%%%%%%%
\subsubsection{Treating the Church--Turing Thesis as a Formal Theorem}

The thesis relates an informal concept of effective calculability to a
formal mathematical notion.

Because one side is informal, it is not proved in the same way as an
ordinary theorem.

%%%%%%%%%%%%%%%%%%%%%%%%%%%%%%%%%%%%%%%%%%%%%%%%%%%%%%%%%%%%
\subsubsection{Assuming Nondeterminism Solves Undecidable Problems}

Nondeterministic Turing machines have the same computability power as
deterministic Turing machines.

They do not decide the Halting Problem.

%%%%%%%%%%%%%%%%%%%%%%%%%%%%%%%%%%%%%%%%%%%%%%%%%%%%%%%%%%%%
\subsubsection{Assuming More Tapes Change Computability}

Multitape machines may improve efficiency.

They do not enlarge the class of computable functions.

%%%%%%%%%%%%%%%%%%%%%%%%%%%%%%%%%%%%%%%%%%%%%%%%%%%%%%%%%%%%
\subsubsection{Using the Wrong Reduction Direction}

To prove target \(B\) undecidable from known-undecidable \(A\), establish

\[
\boxed{
A\leq_mB.
}
\]

Reversing the direction does not prove the desired result.

%%%%%%%%%%%%%%%%%%%%%%%%%%%%%%%%%%%%%%%%%%%%%%%%%%%%%%%%%%%%
\subsubsection{Failing to Prove Both Directions of a Mapping Reduction}

One must show

\[
\boxed{
w\in A
\iff
f(w)\in B.
}
\]

Showing only one implication is not enough.

%%%%%%%%%%%%%%%%%%%%%%%%%%%%%%%%%%%%%%%%%%%%%%%%%%%%%%%%%%%%
\subsubsection{Using an Uncomputable Reduction}

A valid mapping reduction requires the transformation function \(f\) itself
to be computable.

An abstract mathematical correspondence is insufficient.

%%%%%%%%%%%%%%%%%%%%%%%%%%%%%%%%%%%%%%%%%%%%%%%%%%%%%%%%%%%%
\subsubsection{Applying Rice's Theorem to Syntactic Properties}

Rice's theorem concerns nontrivial semantic properties of recognized
languages or computed functions.

It does not apply to simple properties of program text.

%%%%%%%%%%%%%%%%%%%%%%%%%%%%%%%%%%%%%%%%%%%%%%%%%%%%%%%%%%%%
\subsubsection{Applying Rice's Theorem Without Checking Nontriviality}

A property true for every recognizable language or for none of them is
trivial.

Rice's theorem requires a nontrivial property.

%%%%%%%%%%%%%%%%%%%%%%%%%%%%%%%%%%%%%%%%%%%%%%%%%%%%%%%%%%%%
\subsubsection{Assuming Halting Is Not Recognizable}

Halting is recognizable.

The impossibility is deciding both halting and nonhalting in every case.

%%%%%%%%%%%%%%%%%%%%%%%%%%%%%%%%%%%%%%%%%%%%%%%%%%%%%%%%%%%%
\subsubsection{Assuming the Complement of Every Recognizable Language Is Recognizable}

Recognizable languages are not closed under complement.

If both a language and its complement are recognizable, then the language
is decidable.

%%%%%%%%%%%%%%%%%%%%%%%%%%%%%%%%%%%%%%%%%%%%%%%%%%%%%%%%%%%%
\subsubsection{Confusing Finite Bounded Simulation With the Halting Problem}

Given a fixed time bound \(t\), simulation for \(t\) steps is decidable.

The undecidable problem asks whether a machine halts after any finite number
of steps, with no computable bound supplied.

%%%%%%%%%%%%%%%%%%%%%%%%%%%%%%%%%%%%%%%%%%%%%%%%%%%%%%%%%%%%
\subsubsection{Assuming Faster Hardware Can Solve Undecidable Problems}

Undecidability is a mathematical impossibility result for algorithms.

It is not a statement about current hardware limitations.

%%%%%%%%%%%%%%%%%%%%%%%%%%%%%%%%%%%%%%%%%%%%%%%%%%%%%%%%%%%%
\subsubsection{Confusing Mapping Reducibility and Turing Reducibility}

A mapping reduction uses one computable transformation.

A Turing reduction may make repeated adaptive oracle queries.

Turing reducibility is more general.

%%%%%%%%%%%%%%%%%%%%%%%%%%%%%%%%%%%%%%%%%%%%%%%%%%%%%%%%%%%%
\subsubsection{Confusing Gödel Incompleteness With the Halting Problem}

The results are deeply related but distinct.

Gödel incompleteness concerns provability in formal arithmetic theories.

The Halting Problem concerns algorithmic termination.

Both exploit encoding and self-reference.

%%%%%%%%%%%%%%%%%%%%%%%%%%%%%%%%%%%%%%%%%%%%%%%%%%%%%%%%%%%%
\subsection{Applications}
\label{subsec:computability-applications}
%%%%%%%%%%%%%%%%%%%%%%%%%%%%%%%%%%%%%%%%%%%%%%%%%%%%%%%%%%%%

\begin{applicationbox}

\textbf{Foundations of Mathematics.}
Computability gives a precise framework for distinguishing algorithmically
solvable and unsolvable mathematical problems.

\medskip

\textbf{Logic.}
Decision procedures, theorem sets, effective axiomatizations, and
undecidable theories connect computation with formal logic.

\medskip

\textbf{Proof Theory.}
Gödel coding, computable proof checking, and enumeration of proofs connect
formal derivability with effective computation.

\medskip

\textbf{Computer Science.}
Halting, program equivalence, verification limits, and universal
computation form the theoretical foundation of programming-language and
algorithmic limitations.

\medskip

\textbf{Formal Verification.}
Undecidability results explain why general-purpose verification tools must
use restrictions, approximations, annotations, or incomplete procedures.

\medskip

\textbf{Recursion Theory.}
Computable functions, computably enumerable sets, Turing reducibility, and
jumps form the starting point for finer analysis of degrees of
unsolvability.

\medskip

\textbf{Complexity Theory.}
Computability determines whether an algorithm exists at all; complexity
then studies the resources required by such algorithms.

\end{applicationbox}

%%%%%%%%%%%%%%%%%%%%%%%%%%%%%%%%%%%%%%%%%%%%%%%%%%%%%%%%%%%%
\subsection{Exercises}
\label{subsec:computability-exercises}
%%%%%%%%%%%%%%%%%%%%%%%%%%%%%%%%%%%%%%%%%%%%%%%%%%%%%%%%%%%%

\begin{exercise}[1]
Define a decision problem.
\end{exercise}

\begin{exercise}[1]
Define a formal language.
\end{exercise}

\begin{exercise}[1]
Define a Turing machine conceptually.
\end{exercise}

\begin{exercise}[1]
Define a configuration.
\end{exercise}

\begin{exercise}[1]
Define a recognizer.
\end{exercise}

\begin{exercise}[1]
Define a decider.
\end{exercise}

\begin{exercise}[1]
Define a computable function.
\end{exercise}

\begin{exercise}[1]
Define a partial computable function.
\end{exercise}

\begin{exercise}[1]
Define a computably enumerable set.
\end{exercise}

\begin{exercise}[1]
State the Church--Turing thesis.
\end{exercise}

\begin{exercise}[2]
Describe the components of a standard Turing machine.
\end{exercise}

\begin{exercise}[2]
Write a transition rule and explain each component.
\end{exercise}

\begin{exercise}[2]
Describe a machine that decides whether a binary string ends in \(0\).
\end{exercise}

\begin{exercise}[2]
Explain the difference between halting in reject and looping forever.
\end{exercise}

\begin{exercise}[2]
Explain why every decidable language is recognizable.
\end{exercise}

\begin{exercise}[3]
Give an example of a recognizable but undecidable language.
\end{exercise}

\begin{exercise}[3]
Explain why a characteristic function must be total to decide a set.
\end{exercise}

\begin{exercise}[2]
Define an enumerator.
\end{exercise}

\begin{exercise}[3]
Explain why every recognizable language can be enumerated.
\end{exercise}

\begin{exercise}[3]
Explain why dovetailing is necessary in that construction.
\end{exercise}

\begin{exercise}[3]
Design a dovetailing schedule for four computations.
\end{exercise}

\begin{exercise}[2]
Explain why multitape machines do not increase computability power.
\end{exercise}

\begin{exercise}[2]
Explain why nondeterministic machines do not increase computability power.
\end{exercise}

\begin{exercise}[3]
Explain why the equivalence of many computational models supports the
Church--Turing thesis.
\end{exercise}

\begin{exercise}[2]
Explain how machines can be encoded as strings.
\end{exercise}

\begin{exercise}[2]
Explain why there are only countably many Turing machines.
\end{exercise}

\begin{exercise}[3]
Show that there are uncountably many functions

\[
f:\mathbb{N}\to\{0,1\}.
\]
\end{exercise}

\begin{exercise}[3]
Use the previous results to show that noncomputable functions exist.
\end{exercise}

\begin{exercise}[2]
Define a universal Turing machine.
\end{exercise}

\begin{exercise}[3]
Explain why universal computation is related to stored-program computers.
\end{exercise}

\begin{exercise}[2]
Define

\[
A_{\mathrm{TM}}.
\]
\end{exercise}

\begin{exercise}[2]
Explain why

\[
A_{\mathrm{TM}}
\]

is recognizable.
\end{exercise}

\begin{exercise}[2]
Define

\[
\operatorname{HALT}_{\mathrm{TM}}.
\]
\end{exercise}

\begin{exercise}[3]
Explain why the Halting Problem is recognizable.
\end{exercise}

\begin{exercise}[3]
Reproduce the diagonal proof that the Halting Problem is undecidable.
\end{exercise}

\begin{exercise}[3]
Explain the role of self-application in the proof.
\end{exercise}

\begin{exercise}[3]
Compare the Halting Problem proof with Cantor diagonalization.
\end{exercise}

\begin{exercise}[2]
Define the complement of a language.
\end{exercise}

\begin{exercise}[2]
Prove decidable languages are closed under complement.
\end{exercise}

\begin{exercise}[3]
Prove that a language is decidable if both it and its complement are
recognizable.
\end{exercise}

\begin{exercise}[3]
Use this result to show that the complement of the Halting Problem is not
recognizable.
\end{exercise}

\begin{exercise}[3]
Prove decidable languages are closed under union.
\end{exercise}

\begin{exercise}[3]
Prove decidable languages are closed under intersection.
\end{exercise}

\begin{exercise}[3]
Show recognizable languages are closed under union.
\end{exercise}

\begin{exercise}[3]
Show recognizable languages are closed under intersection.
\end{exercise}

\begin{exercise}[2]
Define a mapping reduction.
\end{exercise}

\begin{exercise}[3]
Explain why a mapping reduction must be computable.
\end{exercise}

\begin{exercise}[3]
Explain the correct reduction direction for proving undecidability.
\end{exercise}

\begin{exercise}[3]
Prove that if

\[
A\leq_mB
\]

and \(B\) is decidable, then \(A\) is decidable.
\end{exercise}

\begin{exercise}[3]
Deduce that if \(A\) is undecidable and

\[
A\leq_mB,
\]

then \(B\) is undecidable.
\end{exercise}

\begin{exercise}[3]
Explain why both directions of

\[
w\in A
\iff
f(w)\in B
\]

are required.
\end{exercise}

\begin{exercise}[2]
Define the Turing-machine emptiness problem.
\end{exercise}

\begin{exercise}[3]
Explain conceptually why the emptiness problem is undecidable.
\end{exercise}

\begin{exercise}[2]
Define the Turing-machine equivalence problem.
\end{exercise}

\begin{exercise}[3]
Explain why general program equivalence is undecidable.
\end{exercise}

\begin{exercise}[2]
State Rice's theorem.
\end{exercise}

\begin{exercise}[3]
Explain the meaning of a nontrivial semantic property.
\end{exercise}

\begin{exercise}[3]
Give three properties to which Rice's theorem applies.
\end{exercise}

\begin{exercise}[3]
Give a syntactic machine property to which Rice's theorem does not apply.
\end{exercise}

\begin{exercise}[3]
Use Rice's theorem to show that determining whether

\[
L(M)=\varnothing
\]

is undecidable.
\end{exercise}

\begin{exercise}[3]
Use Rice's theorem to analyze whether a machine recognizes a finite
language.
\end{exercise}

\begin{exercise}[2]
State the Post Correspondence Problem.
\end{exercise}

\begin{exercise}[2]
Find a solution to a small supplied PCP instance.
\end{exercise}

\begin{exercise}[3]
Explain why PCP is useful in undecidability reductions.
\end{exercise}

\begin{exercise}[2]
Explain why bounded halting is decidable.
\end{exercise}

\begin{exercise}[3]
Distinguish bounded simulation from the ordinary Halting Problem.
\end{exercise}

\begin{exercise}[3]
Explain why unbounded witness search may fail to decide nonmembership.
\end{exercise}

\begin{exercise}[3]
Show that sets of the form

\[
n\in A
\iff
\exists y\,R(n,y),
\]

with computable \(R\), are computably enumerable.
\end{exercise}

\begin{exercise}[3]
Explain the difference between countable and effectively enumerable.
\end{exercise}

\begin{exercise}[2]
Describe a computable pairing function conceptually.
\end{exercise}

\begin{exercise}[3]
Explain why finite sequences can be encoded by natural numbers.
\end{exercise}

\begin{exercise}[3]
Explain why a finite halting computation can be mechanically verified.
\end{exercise}

\begin{exercise}[3]
Explain universal search.
\end{exercise}

\begin{exercise}[2]
Distinguish computability and computational complexity.
\end{exercise}

\begin{exercise}[3]
Give an example of a decidable problem solvable by exponential brute-force
search.
\end{exercise}

\begin{exercise}[3]
Explain why undecidability is stronger than intractability.
\end{exercise}

\begin{exercise}[2]
Define oracle computation.
\end{exercise}

\begin{exercise}[3]
Define relative computability.
\end{exercise}

\begin{exercise}[3]
Define Turing reducibility.
\end{exercise}

\begin{exercise}[3]
Compare mapping reducibility and Turing reducibility.
\end{exercise}

\begin{exercise}[3]
Define a Turing degree conceptually.
\end{exercise}

\begin{exercise}[3]
Explain the role of a halting oracle.
\end{exercise}

\begin{exercise}[3]
Explain the Turing jump conceptually.
\end{exercise}

\begin{exercise}[3]
Explain why the jump represents stronger computational information.
\end{exercise}

\begin{exercise}[3]
Describe the arithmetic hierarchy conceptually.
\end{exercise}

\begin{exercise}[3]
Explain why existential computable-witness properties are associated with

\[
\Sigma^0_1.
\]
\end{exercise}

\begin{exercise}[3]
Explain why universal computable conditions resemble

\[
\Pi^0_1
\]

properties.
\end{exercise}

\begin{exercise}[3]
Explain why the Halting Problem does not prevent proving termination of
specific programs.
\end{exercise}

\begin{exercise}[3]
Construct a decreasing measure proving termination of a simple recursive
algorithm.
\end{exercise}

\begin{exercise}[3]
Explain why the theorem set of an effectively axiomatized proof system is
computably enumerable.
\end{exercise}

\begin{exercise}[3]
Distinguish a decidable theory from an effectively axiomatized theory.
\end{exercise}

\begin{exercise}[3]
Explain the connection between computability and Gödel coding.
\end{exercise}

\begin{exercise}[3]
Explain how self-reference appears in both the Halting Problem and
incompleteness.
\end{exercise}

\begin{exercise}[4]
Construct a formal Turing machine for binary increment.
\end{exercise}

\begin{exercise}[4]
Construct a Turing machine recognizing

\[
\{0^n1^n:n\geq0\}.
\]
\end{exercise}

\begin{exercise}[4]
Study simulation of multitape Turing machines by single-tape machines.
\end{exercise}

\begin{exercise}[4]
Study simulation of nondeterministic Turing machines by deterministic
machines.
\end{exercise}

\begin{exercise}[4]
Prove equivalence between Turing recognizability and effective
enumerability.
\end{exercise}

\begin{exercise}[4]
Study universal Turing machines in detail.
\end{exercise}

\begin{exercise}[4]
Study effective machine encodings.
\end{exercise}

\begin{exercise}[4]
Prove undecidability of

\[
A_{\mathrm{TM}}.
\]
\end{exercise}

\begin{exercise}[4]
Prove undecidability of

\[
E_{\mathrm{TM}}.
\]
\end{exercise}

\begin{exercise}[4]
Prove undecidability of

\[
EQ_{\mathrm{TM}}.
\]
\end{exercise}

\begin{exercise}[4]
Study the proof of Rice's theorem.
\end{exercise}

\begin{exercise}[4]
Study the Post Correspondence Problem reduction framework.
\end{exercise}

\begin{exercise}[4]
Study many-one degrees.
\end{exercise}

\begin{exercise}[4]
Study Turing degrees.
\end{exercise}

\begin{exercise}[4]
Study the Turing jump.
\end{exercise}

\begin{exercise}[4]
Study the arithmetic hierarchy.
\end{exercise}

\begin{exercise}[4]
Study computable inseparability.
\end{exercise}

\begin{exercise}[4]
Study productive and creative sets.
\end{exercise}

\begin{exercise}[4]
Study the recursion theorem.
\end{exercise}

\begin{exercise}[4]
Study decidable first-order theories.
\end{exercise}

\begin{exercise}[4]
Study undecidability of sufficiently expressive arithmetic theories.
\end{exercise}

%%%%%%%%%%%%%%%%%%%%%%%%%%%%%%%%%%%%%%%%%%%%%%%%%%%%%%%%%%%%
\subsection{Conceptual Summary}
\label{subsec:computability-summary}
%%%%%%%%%%%%%%%%%%%%%%%%%%%%%%%%%%%%%%%%%%%%%%%%%%%%%%%%%%%%

Computability theory asks

\[
\boxed{
\text{which problems can be solved algorithmically}.
}
\]

A Turing machine formalizes effective computation.

A language is decidable when some machine always halts and correctly
determines membership.

A language is recognizable when positive instances are eventually
accepted, while negative instances may cause nontermination.

Thus

\[
\boxed{
\text{decidable}
\Longrightarrow
\text{recognizable},
}
\]

but not conversely.

The Halting Problem provides the central example:

\[
\boxed{
\operatorname{HALT}_{\mathrm{TM}}
\text{ is recognizable but undecidable}.
}
\]

Reductions transfer undecidability:

\[
\boxed{
A\leq_mB,
\qquad
A
\text{ undecidable}
\Longrightarrow
B
\text{ undecidable}.
}
\]

Rice's theorem shows that nontrivial semantic properties of arbitrary
Turing-machine behavior are undecidable.

Universal computation permits machines to manipulate descriptions of other
machines, creating the possibility of self-reference.

Relative computability extends the theory using oracle access:

\[
\boxed{
A\leq_TB.
}
\]

Computability theory therefore identifies fundamental mathematical limits
on algorithms while providing the foundation for recursion theory,
complexity theory, proof theory, and formal computer science.

%%%%%%%%%%%%%%%%%%%%%%%%%%%%%%%%%%%%%%%%%%%%%%%%%%%%%%%%%%%%
\subsection{Section Connections}
\label{subsec:computability-connections}
%%%%%%%%%%%%%%%%%%%%%%%%%%%%%%%%%%%%%%%%%%%%%%%%%%%%%%%%%%%%

\chapterconnection{
Computability Theory develops the algorithmic side of the logical framework
introduced earlier in Chapter~12. Formal Logic specifies languages and
derivations, Model Theory studies semantic structures, and Proof Theory
treats proofs as formal symbolic objects. Computability Theory asks which
such symbolic processes can be carried out algorithmically and establishes
fundamental limitations through diagonalization, the Halting Problem,
reductions, Rice's theorem, and relative computability. The next section,
Recursion Theory, develops computable and computably enumerable sets at a
deeper structural level through recursive-function definitions, indices,
the recursion theorem, reducibility, degrees, jumps, and the arithmetic
hierarchy.
}

%%%%%%%%%%%%%%%%%%%%%%%%%%%%%%%%%%%%%%%%%%%%%%%%%%%%%%%%%%%%
% SECTION SOURCE TRACKING
%%%%%%%%%%%%%%%%%%%%%%%%%%%%%%%%%%%%%%%%%%%%%%%%%%%%%%%%%%%%

%------------------------------------------------------------
% 12.5 Computability Theory
%------------------------------------------------------------
%
% Source basis:
%   - Standard classical computability theory.
%   - Standard Turing-machine theory.
%   - Standard undecidability, reducibility, and relative-computability
%     concepts.
%
% Used for:
%   - decision problems;
%   - languages;
%   - effective procedures;
%   - Turing machines;
%   - tapes, states, symbols, heads, and transitions;
%   - configurations;
%   - halting;
%   - recognizers and deciders;
%   - computable and partial computable functions;
%   - characteristic functions;
%   - computably enumerable sets;
%   - enumerators;
%   - dovetailing;
%   - equivalent computation models;
%   - Church--Turing thesis;
%   - nondeterministic and multitape machines;
%   - encoding machines and inputs;
%   - countability of programs;
%   - existence of noncomputable functions;
%   - universal Turing machines;
%   - universal computable functions;
%   - acceptance problem;
%   - Halting Problem;
%   - diagonalization;
%   - recognizable but undecidable sets;
%   - complements;
%   - recognizable and co-recognizable sets;
%   - closure properties;
%   - mapping reductions;
%   - reduction direction;
%   - emptiness and equivalence problems;
%   - Rice's theorem;
%   - syntactic versus semantic properties;
%   - Post Correspondence Problem;
%   - bounded simulation;
%   - finite versus infinite search;
%   - witness search;
%   - effective enumeration;
%   - pairing and coding;
%   - computation histories;
%   - universal search;
%   - computability versus complexity;
%   - oracle computation;
%   - relative computability;
%   - Turing reducibility;
%   - Turing degrees;
%   - Turing jump;
%   - arithmetic hierarchy preview;
%   - program verification limits;
%   - termination proofs;
%   - theoremhood and effective enumeration;
%   - decidable theories;
%   - self-reference;
%   - fixed-point phenomena preview;
%   - applications and exercises.
%
% Notes:
%   - Recursion Theory is developed next in Section 12.6.
%   - Proof Theory appears in Section 12.4.
%   - Type Theory follows in Section 12.7.
%   - Constructive Mathematics follows in Section 12.8.
%
%%%%%%%%%%%%%%%%%%%%%%%%%%%%%%%%%%%%%%%%%%%%%%%%%%%%%%%%%%%%